\section{Decision Windows Revisited}

This biological layer sharpens the meaning of a decision window. A decision
window is not merely a motivational slogan. It can be read as a brief interval
in which the current regime has weakened enough that a small intervention can
redirect the trajectory before the previous basin regains control.

\begin{discussion}
This does not replace Chapter~\ref{ch:dynamics}. It deepens it. The earlier
chapter identified windows behaviourally as low-barrier intervals. The present
chapter adds a biological reading: windows are credible because organized
systems can approach loss of stability, become more sensitive to small
perturbations, and then settle into a different regime.
\end{discussion}

\begin{example}[Sleep onset as a transition model]
The sleep-transition source is useful precisely because it is not a metaphor
invented by the manuscript. The indexed Nature Neuroscience article treats
falling asleep as a predictable bifurcation dynamic. That does not mean every
behavioural switch is literally the same mechanism, but it does justify asking
a sharper question in daily life: which control parameter is being moved, how
close is the current regime to losing stability, and how large an intervention
is needed before the path really bends?
\end{example}

\begin{example}[Behavioural reinterpretation]
Suppose an agent returns home intending to study. If fatigue, device cues, and
setup friction still stabilize the old leisure regime, the window is weak and
closes quickly. If the desk is prepared, the phone is absent, and the first
exercise is already visible, a small intervention may be enough to redirect
the path before the old basin reasserts itself. The chapter does not claim
that this is literally the sleep transition mechanism; it claims that the
language of regime vulnerability is now better motivated.
\end{example}
